%%
%% spring-lecture08.tex
%% 
%% Made by Alex Nelson <pqnelson@gmail.com>
%% Login   <alex@lisp>
%% 
%% Started on  2026-04-16T08:46:53-0700
%% Last update 2026-04-16T08:46:53-0700
%% 

\lecture{}

\begin{node}
We will be sloppy with our notation, writing $g(X,Y)$ like
differential geometers do, instead of $g[X,Y]$ like analysts do.
\end{node}

\begin{lemma}
Let $M$ be a smooth $n$-manifold. Let
$(U,\varphi=(\xi^{1},\dots,\xi^{n}))\in\smoothstr{M}$ be a local chart
of $M$. Then the coordinate vector fields $\partial/\partial\xi^{i}$
have a vanishing Lie bracket with each other:
\begin{equation}
\left[\frac{\partial}{\partial\xi^{i}},\frac{\partial}{\partial\xi^{j}}\right]=0
\end{equation}
for all $i,j=1,\dots,n$.
\end{lemma}

\begin{proof}
Let $f\in C^{\infty}(M)$ be any smooth function on $M$. Then we see locally
\begin{equation}
\left[\frac{\partial}{\partial\xi^{i}},\frac{\partial}{\partial\xi^{j}}\right](f)=
\frac{\partial}{\partial\xi^{i}}\left(\frac{\partial f}{\partial\xi^{j}}\right)
-\frac{\partial}{\partial\xi^{j}}\left(\frac{\partial f}{\partial\xi^{i}}\right).
\end{equation}
But we see
\begin{equation}
\frac{\partial}{\partial\xi^{i}}\left(\frac{\partial f}{\partial\xi^{j}}\right)
=\frac{\partial}{\partial\xi^{i}}\left(\frac{\partial(f\circ\varphi^{-1})}{\partial\xi^{j}}\circ\varphi\right).
\end{equation}
Let us define
\begin{equation}
g:=\frac{\partial(f\circ\varphi^{-1})}{\partial\xi^{j}}\circ\varphi
\end{equation}
\begin{subequations}
  \begin{align}
\frac{\partial}{\partial\xi^{i}}\left(\frac{\partial f}{\partial\xi^{j}}\right)
&=\frac{\partial}{\partial\xi^{i}}g\\
&=\left(\frac{\partial}{\partial\xi^{i}}[g\circ\varphi^{-1}]\right)\circ\varphi\\
&=\left(\frac{\partial}{\partial\xi^{i}}\left[\left(\frac{\partial(f\circ\varphi^{-1})}{\partial\xi^{j}}\right)\circ\varphi^{-1}\circ\varphi\right]\right)\circ\varphi\\
&=\left(\frac{\partial}{\partial\xi^{i}}\left[\frac{\partial(f\circ\varphi^{-1})}{\partial\xi^{j}}\right]\right)\circ\varphi\\
&=\left(\frac{\partial}{\partial\xi^{j}}\left[\frac{\partial(f\circ\varphi^{-1})}{\partial\xi^{i}}\right]\right)\circ\varphi\quad\mbox{by Scharz's theorem}\\
&=\dots\quad\mbox{same reasoning unrolled backwards}\\
&=\frac{\partial}{\partial\xi^{j}}\left(\frac{\partial f}{\partial\xi^{i}}\right).
  \end{align}
\end{subequations}
Hence the Lie bracket vanishes.
\end{proof}

\begin{proposition}
Let $(M,g)$ be a Riemannian manifold. Let $(U,\varphi=(\xi^{1},\dots,\xi^{n}))\in\smoothstr{M}$
be a local chart on $M$. Then the Christoffel symbols for the
Levi-Civita connection compatible with $g$ in $(U,\varphi)$ are given by:
\begin{equation}
\Gamma^{k}_{ij}=\frac{1}{2}\sum^{n}_{\ell=1}g^{k\ell}\left(\frac{\partial g_{j\ell}}{\partial\xi^{i}}
  +\frac{\partial g_{\ell i}}{\partial\xi^{j}}
  -\frac{\partial g_{ij}}{\partial\xi^{\ell}}\right)
\end{equation}
for all $i,j,k=1,\dots,n$ where $(g^{ij})$ is the matrix inverse of
the matrix $(g_{ij})$: $\sum^{n}_{\ell=1}g^{i\ell}g_{\ell j}=\delta^{i}_{j}$
for every $i,j=1,\dots,n$.
\end{proposition}

\begin{proof}
Note that the Lie bracket of two coordinate vector fields vanish:
\begin{equation}
\left[\frac{\partial}{\partial\xi^{i}},\frac{\partial}{\partial\xi^{j}}\right]=0,
\end{equation}
for all $i,j=1,\dots,n$. Then applying this to Theorem~\ref{thm:fundamental:riemannian-geometry},
\begin{equation}
g(\nabla_{\partial/\partial\xi^{i}}\frac{\partial}{\partial\xi^{j}},\frac{\partial}{\partial\xi^{\ell}})=\frac{1}{2}\left(\frac{\partial g_{j\ell}}{\partial\xi^{i}}
+\frac{\partial g_{i\ell}}{\partial\xi^{j}}-\frac{\partial g_{ij}}{\partial\xi^{\ell}}\right).
\end{equation}
But the $\partial/\partial\xi^{i}$ are not necessarily an orthonormal
basis for vector fields, so we need
\begin{subequations}
  \begin{align}
\nabla_{\partial/\partial\xi^{i}}\frac{\partial}{\partial\xi^{j}}
&=\sum^{n}_{k=1}\left(\sum^{n}_{\ell=1}g^{k\ell}g(\nabla_{\partial/\partial\xi^{i}}\frac{\partial}{\partial\xi^{j}},\frac{\partial}{\partial\xi^{\ell}}\right)\frac{\partial}{\partial\xi^{k}}
\quad\mbox{by linear algebra}\\
&=\sum^{n}_{k=1}\left(\frac{1}{2}\sum^{n}_{\ell=1}%
g^{k\ell}\left(\frac{\partial g_{i\ell}}{\partial\xi^{j}}
  +\frac{\partial g_{j\ell}}{\partial\xi^{i}}
  -\frac{\partial g_{ij}}{\partial\xi^{\ell}}\right)\right)%
  \frac{\partial}{\partial\xi^{k}}\\
&=\sum^{n}_{k=1}\Gamma^{k}_{ij}\frac{\partial}{\partial\xi^{k}},
  \end{align}
\end{subequations}
which implies the claim.
\end{proof}

\begin{remark}
If $V$ is a vector space with basis $e_{1}$, \dots, $e_{n}$ and an
inner product $\langle-,-\rangle$, then we can consider the matrix
\begin{equation}
G=(G)_{ij}=\langle e_{i},e_{j}\rangle
\end{equation}
with matrix inverse $G^{-1}=(g^{ij})$. Then for any $v\in V$, we can
express
\begin{equation}
v=\sum^{n}_{k=1}\left(\sum^{n}_{\ell=1}g^{k\ell}\langle v,e_{\ell}\rangle\right)e_{k}.
\end{equation}
This is precisely what we're doing for $V=T_{x}M$ with inner product
$\langle-,-\rangle=g(x)$ and basis $e_{i}=\partial/\partial\xi^{i}$.
\end{remark}

\begin{remark}
If you have any connection (not necessarily the Levi-Civita
connection) $\nabla$ on a smooth manifold $M$, then it extends
naturally to a connection on the tensor bundle $T^{p}_{q}M$ for every
$p,q\in\NN_{0}$. The construction consists of two steps:
\begin{enumerate}
\item\textsc{Extension to $T^{*}M$}: For any $\omega\in T^{*}M=\Omega^{1}(M)$,
for any $X\in\Vect{M}$, define
$(\nabla_{X}\omega)(Y):=X(\omega(Y))-\omega(\nabla_{X}Y)$ for all $Y\in\Vect{M}$.
Exercise: check this is a connection on $T^{*}M$.
\item\textsc{Extension to tensor products}: For any
  $T\in\Gamma(T^{p}_{q}M)$ and $S\in\Gamma(T^{a}_{b}M)$, for any $X\in\Vect{M}$
  define
  \begin{equation}
\nabla_{X}(T\otimes S)=(\nabla_{X}T)\otimes S+T\otimes(\nabla_{X}S).
  \end{equation}
  Not every tensor is a tensor product of vector fields and one-forms
  (they are called \define{Simple Tensors}). Locally, we can do this,
  but not globally.
\end{enumerate}
\end{remark}

\begin{notation}
Let $T\in\Gamma(T^{p}_{q}M)$ be a tensor, then the map $\nabla T$
is given by
\begin{equation}
\nabla T(\omega_{1},\dots,\omega_{p},X,Y_{1},\dots,Y_{q})=(\nabla_{X}T)(\omega_{1},\dots,\omega_{p},Y_{1},\dots,Y_{q}),
\end{equation}
which is a $(p,q+1)$ tensor on $M$.
\end{notation}

\begin{xca}
Check that $\nabla T$ is $C^{\infty}(M)$-linear in all its arguments.
\end{xca}

\begin{theorem}
Let $(M,g)$ be a Riemannian manifold, let $\nabla$ be the Levi-Civita
connection on $M$. Then
\begin{equation}
\nabla g=0,
\end{equation}
i.e., for any $X\in\Vect{M}$ we have $\nabla_{X}g=0$.
\end{theorem}

We only require $\nabla$ to be a metric-compatible connection, not
necessarily torsion-free. \emph{Every} metric-compatible connection
satisfies this condition.

\begin{proof}
Choose any local chart $(U,\varphi=(\xi^{1},\dots,\xi^{n}))\in\smoothstr{M}$
on $M$ and note that on this chart
\begin{equation}
g=\sum^{n}_{i,j=1}g_{i,j}\,\D\xi^{i}\otimes\D\xi^{j}.
\end{equation}
Let us note also that
\begin{equation}
X=\sum^{n}_{k=1}X^{k}\frac{\partial}{\partial\xi^{k}}.
\end{equation}
Then in our local chart,
\begin{equation}
\nabla_{X}g=\sum^{n}_{k=1}X^{k}\sum^{n}_{i,j=1}\nabla_{\partial/\partial\xi^{k}}(g_{ij}\,\D\xi^{i}\otimes\D\xi^{j}).
\end{equation}
We just need to prove
\begin{equation}
\nabla_{\partial/\partial\xi^{k}}(g_{ij}\,\D\xi^{i}\otimes\D\xi^{j})=0
\end{equation}
for all $i,j,k$. We find
\begin{subequations}
  \begin{align}
\nabla_{\partial/\partial\xi^{k}}(\D\xi^{i})
&=\sum^{n}_{\ell=1}(\nabla_{\partial/\partial\xi^{k}}(\D\xi^{i}))\left(\frac{\partial}{\partial\xi^{\ell}}\right)\D\xi^{\ell}\\
&=\sum^{n}_{\ell=1}\left(\frac{\partial}{\partial\xi^{k}}\underbrace{\left(\D\xi^{i}\left(\frac{\partial}{\partial\xi^{\ell}}\right)\right)}_{=\delta^{i}_{\ell}}-\D\xi^{i}(\nabla_{\partial/\partial\xi^{i}}\frac{\partial}{\partial\xi^{\ell}})\right)\,\D\xi^{\ell}\\
&=\sum^{n}_{\ell=1}\left(0-\D\xi^{i}(\nabla_{\partial/\partial\xi^{i}}\frac{\partial}{\partial\xi^{\ell}})\right)\,\D\xi^{\ell}\\
&=\sum^{n}_{\ell=1}-\Gamma^{i}_{k\ell}\,\D\xi^{\ell}.
  \end{align}
\end{subequations}
\end{proof}